\chapter{Aplikasi BFS dan DFS}

Penelusuran \textit{Document Object Model} (DOM) dibangun menggunakan antarmuka berbasis web (React dengan TypeScript) dan sistem \textit{backend} (sebagai \textit{middleware/proxy} menggunakan Vite Plugin) untuk melakukan \textit{scraping} HTML. Setelah mendapatkan HTML dari target, seluruh token didalamnya diproses menjadi struktur data DOM \textit{Tree} lalu algoritma \textit{Breadth First Search} (BFS) dan \textit{Depth First Search} (DFS) diaplikasikan untuk menelusuri elemen-elemen sesuai dengan CSS \textit{Selector} yang diberikan pengguna.

\section{Mekanisme Algoritma Penelusuran}

DFS dan BFS diimplementasikan untuk mengunjungi setiap \textit{node} pada DOM Tree. Algoritma didesain agar dapat menerima sebuah fungsi predikat yang bertugas untuk mengevaluasi apakah suatu \textit{node} memenuhi kriteria CSS \textit{Selector} yang diminta.

\subsection{Depth First Search (DFS)}
Depth First Search merupakan algoritma penelusuran secara mendalam (\textit{deep}). DFS menelusuri cabang pohon hingga mencapai \textit{leaf node} sebelum melakukan \textit{backtracking} untuk mengeksplorasi cabang lainnya.

Mekanisme DFS diimplementasikan dengan pendekatan rekursif dan penelusuran \textit{pre-order}. Terdapat dua fungsi utama untuk DFS:
\begin{itemize}
    \item \textbf{\texttt{dfsWalk}}: Fungsi ini digunakan untuk menelusuri seluruh pohon dan mengeksekusi suatu aksi (misalnya mengumpulkan elemen yang cocok dengan \textit{selector}) pada setiap \textit{node} yang dikunjungi. Algoritma memanggil fungsi secara rekursif pada setiap anak (\textit{children}) dari suatu elemen, sembari mencatat kedalaman (\textit{depth}) iterasi saat ini.
    \item \textbf{\texttt{dfs}}: Fungsi ini digunakan untuk menemukan elemen pertama yang memenuhi syarat (fungsi predikat mengembalikan nilai \textit{true}). Jika elemen saat ini memenuhi kriteria, pencarian dihentikan dan elemen tersebut dikembalikan. Jika tidak, pencarian akan dilanjutkan secara rekursif ke anak-anaknya.
\end{itemize}

\subsection{Breadth First Search (BFS)}
Breadth First Search merupakan algoritma penelusuran secara melebar (\textit{wide}). BFS mengeksplorasi seluruh \textit{node} pada tingkat kedalaman (\textit{depth}) yang sama terlebih dahulu sebelum turun ke tingkat kedalaman berikutnya.

Mekanisme BFS diimplementasikan dengan pendekatan iteratif menggunakan struktur data \textit{queue}. Terdapat dua fungsi utama untuk BFS:
\begin{itemize}
    \item \textbf{\texttt{bfsWalk}}: Fungsi penelusuran lengkap yang menginisialisasi sebuah antrean dengan node root di dalamnya. Selama antrean tidak kosong, \textit{node} paling depan akan diambil (\textit{dequeue}), dievaluasi, dan seluruh anak dari \textit{node} tersebut akan dimasukkan ke bagian belakang antrean (\textit{enqueue}).
    \item \textbf{\texttt{bfs}}: Fungsi pencarian yang iteratif menelusuri pohon tingkat demi tingkat. Jika \textit{node} yang sedang di-\textit{dequeue} memenuhi fungsi predikat, algoritma akan langsung mengembalikan \textit{node} tersebut dan proses perulangan (\textit{while loop}) dihentikan.
\end{itemize}

\section{Pemodelan Pohon DOM dan Integrasi Penelusuran}

Sebelum BFS dan DFS dapat digunakan, aplikasi harus merepresentasikan kode HTML sebagai graf berbentuk pohon berakar (\textit{rooted tree}). Proses ini dilakukan melalui \textit{Scraping} \& \textit{Tokenizing}, serta \textit{Tree Building}.

\subsection{Scraping dan Parsing HTML}
Batasan CORS (\textit{Cross-Origin Resource Sharing}) menyebabkan permintaan data dilakukan melalui \textit{backend middleware} menggunakan \texttt{node:http}. Hasil \textit{scraping} berupa string HTML dikembalikan ke \textit{frontend} dan diproses oleh \textit{library html-tokenizer}.  Teks hierarkis diubah menjadi kumpulan token datar (seperti token \textit{open tag}, \textit{close tag}, \textit{text}, dan \textit{comment}).

\subsection{Pembentukan Pohon (Tree Building)}
Kumpulan token kemudian disusun menjadi pohon menggunakan pendekatan \textit{Stack} yang berisikan \texttt{ElmtNode} (elemen tag), \texttt{TextNode} (teks biasa), dan \texttt{CommentNode} (komentar HTML).
Langkah-langkah penyusunan pohon adalah sebagai berikut:
\begin{enumerate}
    \item Inisialisasi \textit{root node} berjenis elemen dengan \textit{tag} \texttt{\#document} dan masukkan ke dalam \textit{stack}.
    \item Setiap iterasi membaca satu token secara linear lalu mengecek apabila token merupakan \textit{open tag}, maka \texttt{ElmtNode} baru dibuat, disematkan sebagai anak (\textit{child}) dari elemen paling atas di \textit{stack}, dan dimasukkan ke dalam \textit{stack} (kecuali jika merupakan tag \textit{self-closing}).
    \item Jika token berupa teks atau komentar, \textit{node} bersangkutan ditambahkan ke dalam \textit{children} milik elemen teratas di \textit{stack} tanpa memodifikasi \textit{stack}.
    \item Jika token berupa \textit{close tag}, pengecekan dilakukan pada \textit{stack} dari atas ke bawah untuk mencari \textit{open tag} yang bersesuaian, lalu membuang elemen-elemen di atasnya sehingga konteks penyusunan (\textit{parent}) kembali naik ke level sebelumnya.
\end{enumerate}

\subsection{Pencocokan CSS Selector}
Setelah \textit{DOM Tree} selesai disusun ke dalam representasi \texttt{Node}, DFS atau BFS dijalankan dengan titik awal pada \textit{root}. Pemrosesan CSS \textit{Selector} diabstraksi menjadi sebuah fungsi kondisional yang menerima parameter sebuah \textit{node}. 

Fungsi DFS (\texttt{dfsWalk}) maupun BFS (\texttt{bfsWalk}) memberikan urutan penyusunan ruang pencarian. Setiap kali algoritma berhenti di suatu \textit{node}, atribut \textit{node} seperti nama \textit{tag} dan koleksi \textit{properties} akan diperiksa oleh predikat \textit{selector}. Elemen yang lolos evaluasi akan disimpan untuk nantinya dirender pada proses visualisasi di sisi \textit{frontend}.